\documentclass[11pt,a4paper]{article}
\makeatletter\def\input@path{{../../templates/}}\makeatother
\usepackage{avaliacao-poo}
\configurarAvaliacao{Gabarito \textemdash{} Checkpoint 1}{Programação Orientada a Objetos}
\begin{document}
\cabecalhoGabarito
\identificacaoInterna{Variante A}

\textbf{Q1 (25):}\quad a = \codigoInline{1}; b = \codigoInline{true}; c = \codigoInline{2}; d = \codigoInline{false}; e = \codigoInline{1}.

\medskip
\textbf{Q2 (25):}\quad a = F; b = V; c = V; d = F; e = F.

\medskip
\textbf{Q3--Q12 (5 cada):}
\begin{tabularx}{\textwidth}{@{}*{5}{>{\bfseries}l X}@{}}
Q3 & C & Q4 & B & Q5 & A & Q6 & D & Q7 & B \\
Q8 & E & Q9 & D & Q10 & C & Q11 & A & Q12 & D
\end{tabularx}

\medskip
\textbf{Traço de execução.} \codigoInline{principal} e \codigoInline{painel} acessam o mesmo objeto. Primeiro são observados andar 1 e porta aberta; depois, pelo outro nome, andar 2 e porta fechada. \codigoInline{servico} é distinto e termina no andar 1 com a porta aberta.

\vspace{18pt}
\identificacaoInterna{Variante B}

\textbf{Q1 (25):}\quad a = \codigoInline{true}; b = \codigoInline{1}; c = \codigoInline{1}; d = \codigoInline{false}; e = \codigoInline{2}.

\medskip
\textbf{Q2 (25):}\quad a = F; b = F; c = F; d = V; e = V.

\medskip
\textbf{Q3--Q12 (5 cada):}
\begin{tabularx}{\textwidth}{@{}*{5}{>{\bfseries}l X}@{}}
Q3 & E & Q4 & C & Q5 & B & Q6 & A & Q7 & D \\
Q8 & C & Q9 & B & Q10 & E & Q11 & C & Q12 & B
\end{tabularx}

\medskip
\textbf{Traço de execução.} \codigoInline{cabine} e \codigoInline{painel} acessam o mesmo objeto. Primeiro são observados andar 1 e porta aberta; depois, andar 1 e porta fechada. \codigoInline{carga} é distinto e termina no andar 2 com a porta aberta.
\end{document}
